\section{国内外研究现状}
\label{intro:related}

多无人机协同观测需要将分散的信息送至接收中心，并在通信资源约束下支持任务判断。与这一过程直接相关的研究包括路由与资源分配、协同观测中的消息选择、面向任务的语义通信，以及消息丢失后的重传控制。以下按照各类方法解决的问题介绍研究进展，最后归纳本文需要进一步处理的不足。

\subsection{无人机网络路由与通信资源分配}

路由首先需要在不断变化的连接中找到可用路径。自组网按需距离矢量路由（Ad Hoc On-Demand Distance Vector，AODV）在需要通信时发现路径，减少无业务期间的维护开销\cite{aodv}；优化链路状态路由（Optimized Link State Routing，OLSR）主动交换拓扑信息，并用多点中继减少控制消息的重复转发\cite{olsr}。两者分别适合按需建立和持续维护路径，但节点移动较快时，都需要处理链路失效和控制信息过期。贪婪周边无状态路由（Greedy Perimeter Stateless Routing，GPSR）利用邻居和目的节点的位置进行局部转发\cite{gpsr}，降低了维护全局拓扑的需求。地理位置能够反映转发进度，却不能单独反映链路质量和中继电量，因而仍需结合通信状态选择路径。

无人机路由在此基础上进一步考虑运动和能量。Lakew 等归纳了飞行自组网中拓扑变化、路径维护与网络寿命之间的关系\cite{intro4}。基于 Q 学习的地理路由方法 QGeo 利用交互结果更新邻居的转发价值\cite{qgeo}；基于 Q 学习的拓扑感知路由（Q-learning-based Topology-Aware Routing，QTAR）将拓扑状态纳入下一跳选择\cite{qtar}。这类方法通过历史通信结果适应环境变化，其选择偏好由状态和回报共同决定。Li 等提出基于强化学习的节能快速路由，结合电池、信道及一跳邻居状态选择转发节点和功率，并同时考虑投递结果、时延与能耗\cite{drlfr}。Jin 等则面向城市应急网络研究深度强化学习路由，提高路径对网络变化的适应能力\cite{drlrr}。这些工作说明，能量和时延可以直接参与路由决策，而不只是事后统计指标。

路径确定后，资源分配影响链路能够提供的速率和服务时间。Mozaffari 等从部署、信道和能量效率等方面讨论无人机通信设计，并总结优化与机器学习等求解途径\cite{mozaffari2019tutorial}。对于模型和环境信息较为明确的问题，可以依据速率与功率关系构建优化目标；在难以准确获得状态转移时，可以通过交互学习动作价值。前一类方法便于表达约束，后一类方法可利用运行数据适应变化。两者都需要明确有限电量如何分配到连续到达的通信需求，而不能仅根据当前链路速率决定发送强度。

当业务到达和未来信道未知时，队列方法提供了另一种在线分配依据。Tassiulas 和 Ephremides 用队列积压描述尚未满足的传输需求，研究多跳网络的吞吐量与稳定性\cite{tassiulas1992stability}；Neely 的李雅普诺夫优化方法进一步将长期资源约束转化为当前决策中的代价\cite{neely2010stochastic}。反压路由依据相邻节点的队列差分配转发量，Ying 等将其与最短路径结合以减少绕行\cite{related2}，Yu 和 Neely 分析了联合速率控制与路由中的效用和队列长度\cite{algo1}。在无人机系统中，Lin 等利用队列和能量状态安排太阳能无人机的资源分配与轨迹\cite{intro2}。这些研究表明，队列可用于联系当前负载与长期资源使用，但具体动作和约束仍由应用条件决定。

对本文的连续观测场景，路由与功率控制需要同时响应两个变化：一是当前数据积压和链路速率，二是关键节点的剩余电量及此前消耗。降低一次传输的能耗有助于节能，但持续使用同一中继仍可能缩短服务时间；均匀轮换中继也可能因使用较差链路而降低吞吐量。本文据此将本地队列、链路代价和累计能耗共同用于分布式决策，研究有限通信电量下的转发能力与网络可用时间。

\subsection{协同观测中的消息选择}

多节点协同观测利用不同视角补充单个节点看不到的信息，但节点增加也会带来更多待传数据。消息选择研究的基本问题是：在允许交换的数据量有限时，确定哪些节点和哪些内容最有助于接收端完成判断。无人机和车辆协同感知都存在局部视野不完整、观测重叠和无线带宽有限的问题，因而其信息筛选方法与本文直接相关；具体选择对象则随图像、特征或文本等消息形式而变化。

早期协同感知方法主要研究如何融合其他节点的信息。车车通信网络 V2VNet 交换压缩的中间特征，用于联合感知与预测\cite{v2vnet}；DiscoNet 通过知识蒸馏学习协作关系，对来自不同节点的特征分配权重\cite{disconet}。这些方法说明了局部观测之间的互补价值，但融合更多特征也需要更多传输资源。因此，后续工作将通信对象和发送内容纳入决策，而不再默认所有节点均提供完整特征。

Who2com 通过可学习的请求与匹配机制选择信息提供者\cite{who2com}，When2com 进一步根据观测需求决定通信时机和协作关系\cite{when2com}。Where2comm 利用空间置信图筛选对感知结果重要的稀疏区域，并随可用带宽调整交换范围，其实验同时涉及车辆和无人机协同感知\cite{where2comm}。How2comm 从空间位置和特征通道两个维度减少冗余，并补偿传输延迟造成的特征不同步\cite{how2comm}。这一发展过程体现出明确的变化：通信选择从确定信息来源，逐步细化到来源中具体需要发送的部分。

消息价值还与先前的通信有关。接收端已经获得某一区域的信息后，再发送相似内容的作用可能下降；一个单独使用时不够可靠的视角，也可能补充其他视角缺少的证据。学习式通信调度可以根据任务反馈决定是否继续交换信息\cite{ic3net,schednet}。本文采用中心选择消息的方式，借鉴的是这种根据任务状态决定发送时机的思路，并将接收端已有判断、来源质量和已选历史作为选择依据。

现有空间筛选方法主要评价图像区域或中间特征，而类别判断、陈述和摘要具有不同的长度与信息密度。对这类候选，仅按单条消息的置信度排序，容易忽略其与已有内容的重复程度。本文第二项研究因此同时选择消息来源和表达粒度，并在每次选择后更新候选评价，在消息总长度约束内提高中心判断的准确性。

\subsection{面向任务的语义通信}

消息选择需要回答信息对任务有多大作用。经典通信理论通过熵、信道容量和率失真描述有限资源下的信息传输与恢复\cite{shannon1948,cover2006information}。对于协同判断，接收端通常只需要观测中与当前问题相关的部分。信息瓶颈方法通过限制中间表示保留的输入信息、维持其与目标的相关性，为任务相关的信息压缩提供了分析依据\cite{tishby2000ib}。面向任务的语义通信据此将最终判断或推理结果纳入通信评价\cite{semantic-communications-jsac,zhou2024goal}。

深度联合信源信道编码通过联合训练编码器和解码器，学习适应信道的传输表示。Bourtsoulatze 等研究图像到信道符号的映射，改善受限带宽和噪声条件下的图像传输\cite{bourtsoulatze2019deepjscc}。深度语义通信模型（Deep Semantic Communication，DeepSC）利用注意力网络提取并传输文本含义\cite{xie2021deepsc,vaswani2017attention}；Shao 等面向边缘推理，以信息瓶颈保留有效特征，并采用可变长度表示适应信道\cite{shao2021task}。这些方法说明，保留任务所需信息能够减少对完整数据恢复的依赖。用于多源观测时，还需进一步确定应当编码和发送哪些来源，才能避免对重复内容进行高效但无益的传输。

视觉语言模型（Vision-Language Model，VLM）为构造显式观测消息提供了工具。对比语言图像预训练模型（Contrastive Language-Image Pre-training，CLIP）通过图文学习建立共享表示\cite{clip}；视觉指令模型能够根据图像和问题生成回答或描述\cite{blip2,llava}。局部观测由此可以表示为类别判断、置信度及简短文本，供接收中心融合使用。模型输出仍可能受遮挡和视角偏差影响，因此生成消息后还需评价其来源及其对当前判断的补充作用。

任务查询决定需要保留的信息。视觉问答（Visual Question Answering，VQA）研究根据问题定位相关区域\cite{where-to-look-vqa}；视觉问答数据集第二版（VQAv2）利用同一问题对应不同答案的图像对减少语言先验偏差\cite{vqa-v2}，面向真实场景的组合问答数据集 GQA 则构造多步推理问题\cite{gqa}。常见物体上下文数据集（Common Objects in Context，COCO）的图像与标注可用于形成局部观测\cite{coco}。本文利用图像与任务查询生成候选消息，并保持提取模型参数固定，使不同通信策略面对相同内容，比较内容选择本身对判断结果和通信量的影响。

\subsection{语义消息的可靠传输与重传控制}

消息经过选择后，仍可能因链路丢失而未被接收。ARQ 依据接收确认和超时安排补发，停等和选择重传分别在协议复杂度、并发能力与资源使用方面作出不同取舍\cite{lin1984arq}。混合自动重传请求（Hybrid Automatic Repeat Request，HARQ）进一步结合差错控制编码与多轮重传。对观测任务而言，重传的收益不仅取决于分组是否缺失，还取决于缺失内容是否会影响最终判断；因此，传输可靠性需要与任务结果联系起来。

语义重传研究已经开始利用任务质量决定补充哪些信息。Sheng 等针对协同感知提取重要特征，结合语义错误检测与 HARQ 改善衰落信道下的感知表现\cite{sheng2024harq}。Li 等研究重复语义信号、多轮接收融合以及增加语义冗余的重传方式\cite{li2025semharq}。这些方法通过任务相关内容和多轮融合提高重传效率，说明传输失败后并非只能按原有字节顺序重复发送全部数据。

重传内容确定后，还存在发送时机的问题。确认消息迟到或丢失时，发送端无法直接区分“接收端缺少该内容”和“内容已到达但尚未获知”两种状态。前一种情况可能需要补发，后一种情况则可能因补发消耗本可用于其他消息的资源。Huang 等用信息效用损失描述收发状态不一致对传输决策的影响，并在卫星融合网络中联合安排采样、编码和重传\cite{uoi2024satellite}。这说明任务收益评价需要使用接收端状态，而不仅是待发消息自身的重要性。

部分可观测马尔可夫决策过程（Partially Observable Markov Decision Process，POMDP）用条件概率表示无法直接观察的状态，并结合历史信息选择动作\cite{kaelbling1998pomdp}。对于观测消息，可以利用发送成功概率、发送记录和已返回的确认估计接收集合，再计算新发或重传的预期收益。考虑更多未来动作有助于分析等待的作用，但会增加在线计算量；模型预测控制通过限制预测时域、按新观测重复求解控制这一开销\cite{mayne2000mpc}。本文第三项研究据此比较有限步内的新发、重传、等待和停止，并将剩余任务时间与通信电量纳入决策。

\subsection{现有研究的不足}

上述工作分别改善了网络传输、信息筛选和重传效率。用于通信资源受限的多无人机协同观测时，仍需把这些能力与连续任务中的实际资源支出联系起来，主要存在以下三个问题。

（1）当前传输收益与后续网络可用时间之间需要进一步协调。能量感知和学习式路由已经将剩余电量、时延及投递结果纳入决策\cite{qtar,drlfr,drlrr}，但同一条高收益路径被持续使用后，关键中继的累计消耗可能改变其后续服务能力。只优化瞬时速率，容易集中消耗这些节点；只按剩余电量分流，又可能牺牲有效转发量。本文关注本地业务积压、链路速率与累计能耗共同作用时的下一跳和功率选择，利用长期能量反馈协调当前转发需求与后续通信能力。

（2）单条消息质量不足以确定其在当前已收信息之后的价值。现有语义表示和协同感知方法能够减少冗余数据\cite{xie2021deepsc,where2comm,how2comm}，但不同来源、不同表达长度的显式消息还具有相互补充或重复的关系。高置信度可能来自受遮挡的视角，较长表达也未必补充中心缺少的判断依据；每次选择还会改变剩余预算与后续消息的作用。本文因此在统一的消息长度预算内，根据来源可靠性、中心判断和通信历史更新候选评价，同时学习消息选择与停止时机。

（3）内容重要性与当前重传必要性需要联合判断。语义重传能够优先补充任务相关信息\cite{sheng2024harq,li2025semharq}，但在确认不及时的条件下，一条重要消息也可能已经到达。发送端需要先估计接收状态，再判断立即补发或等待确认是否更有利。此外，重复消息即使提高了模型置信度，也可能不改变最终类别，因而不应仅按置信度增长估计任务收益。本文第三项研究结合接收状态估计、同类别消息组合的收益一致性和等待决策，减少低收益发送，并把当前节省的能量保留给后续任务。

这三个不足对应观测信息回传中的连续环节。路由研究改善可用传输条件，消息选择研究确定有限资源应承载的内容，发送与重传研究根据实际接收情况调整后续通信。本文以这一关系组织方法设计，并通过最终接收信息是否支持按时正确判断，说明各项研究与任务可靠完成之间的联系。
